\section{Propagation across rank-four components}

Fix a fit, deep tower $T$, coefficient ring $\Lambda=c'.O$, comparison
$\iota$, and the space $X$ and ideal $J\subset B=X.B$ in the original
potential modularity hypothesis. Write $A=X.A$, $d=T.d$ and
$s=\#T.\iota$. A residual point $x:A\to k_\Lambda$ may lie on several
minimal components of $A$. Our goal is to show that if $JA$ vanishes on
one component through $x$, it vanishes on every such component. The
new case is a residual pair of type $A$, where the ordinary local flag
has a nonzero $10$ entry at every marked place. All indices start at zero.

% BEGIN C65R1
\paragraph{C65R1. The tangent equations at a type-$A$ split limit.}
For the tangent calculation work over any field $k$ and any frame
$\Phi$ with $\Phi.\mathrm{Ok}(a)$, where $a:O\to k$. Let $(r,n)$
be a residual pair of type $A$. Its off-diagonal representation
entries have weights in $\{-\beta,-\sigma,-\tau\}$, in the notation
of (UX.46), and $b_v=(n_v)_{10}\ne 0$ at every place.
Filter matrices by the integer $m_i-m_j$ in entry $ij$, with
$m=(0,0,1,1)$. Every allowed off-diagonal entry of $r$ has strictly
positive degree. The degree-zero part of $r$ is its diagonal
$D(g)=\operatorname{diag}(\bar\chi_i(g))$, with
$\bar\chi_i=a\circ\Phi.\chi_i$. A lower unitriangular symplectic
flag has degree-zero part
\[
       n_v^0=1+b_v(E_{10}-E_{32}).
\]
Indeed the symplectic adjoint identity forces $(n_v)_{32}=-(n_v)_{10}$.
In the original flag equation, the first column of $n_v$ is an
eigenvector on the decomposition group $G_v$ of character
$\bar\chi_0$. Since $r_{10}=r_{12}=r_{13}=0$, its row-one equation
is $b_v\bar\chi_1=b_v\bar\chi_0$. Thus $\bar\chi_0=\bar\chi_1$ on
$G_v$, including when the residual field has characteristic two.

For any of the source reduction supports $S$, tangent vectors are
pairs $(X(g),L_v)$ in right logarithmic coordinates (UX.48)--(UX.49).
Taking the first nonzero degree of a tangent vector gives a tangent
vector at $(D,n^0)$ with the same support condition. To check this,
a matrix product has, in its lowest possible degree, the product of
the corresponding lowest parts. Symplectic adjoint preserves the
grading since $m_i+m_{3-i}=1$. Taking the prescribed degree therefore
preserves the linearized cocycle, adjoint, vanishing, normalization
and flag equations; a forced zero entry stays zero in every degree.
Filter the tangent space by successively vanishing degrees. Projection
embeds each successive quotient in the matching homogeneous tangent
space at $(D,n^0)$. These images are independent, since their degrees
differ. There are only three possible degrees. This proves that the
rank at the original pair is at most that at the displayed limit,
even when the ambient function spaces are infinite-dimensional.

At the limit put
$Z=X(h)+D(h)L_vD(h)^{-1}-L_v$ for $h\in G_v$.
Upper triangularity of $(n_v^0)^{-1}Zn_v^0$ gives
\[
 Z_{10}=b_v^2Z_{01}+b_v(Z_{00}-Z_{11}),\qquad
 Z_{21}=Z_{30}=0.
\]
For the last equality the $21,20,31$ equations kill the intervening
entries in the $30$ equation. In particular if $X_{00}=X_{11}=0$,
then locally $X_{10}=b_v^2X_{01}$, since $\bar\chi_0=\bar\chi_1$.
The $21$ and $30$ entries of $X$ are local coboundaries with constants
$(L_v)_{21}$ and $(L_v)_{30}$ respectively. These three consequences
hold at every prescribed place.
% END C65R1

% BEGIN C65R2
\paragraph{C65R2. One unrestricted odd cocycle bounds the tangent rank.}
Assume $\Phi.\mathrm{Bounds}(a,d,C)$. Take a reduction row $S$
containing $-\sigma$, or containing both $-\beta,-\tau$ and
different from $Y=\{\pm\beta,\pm\tau\}$. In the numbering
$0,\ldots,8$ of (UX.46), the following table lists all such rows.
The second entry is one coordinate whose tangent cocycle will be
retained in the quotient.
\[
\begin{array}{c|ccccc}
 S & 0 & 3 & 4 & 6 & 7 \\ \hline
 ij&12&03&01&12&03
\end{array}
\]
In fact $03$ has weight $\tau$, $12$ has weight $\beta$, and
$01,10$ have weights $\alpha,-\alpha$. After the selected entry is
set to zero, the support forces both $03,12$ to vanish globally and
at least one of $01,10$ to vanish globally. Here the tangent support
tests $X(g)D(g)$, which is equivalent entry by entry to testing
$X(g)$ because every diagonal character value is a unit.

At $(D,n^0)$ every $X_{ij}$ is a normalized global cocycle of
weight $w_{ij}$. Project to the selected odd cocycle and to the
two diagonal cocycles $X_{00},X_{11}$. Their total image rank is
at most $(d+C)+2C$ by the original bounds on odd and even weights.
In the kernel, C65R1 forces \emph{both} $01,10$ to vanish on $G_v$,
since $b_v\ne0$. Of the independent odd entries
$01,03,21,12,10,30$, the first, second, fourth and fifth now have
zero local class, while $21,30$ are local coboundaries by C65R1.
Thus every independent odd entry belongs to the strict cocycle space.
All other independent weights are $0,\sigma,-\sigma$. Each strict
odd space and each of these even global spaces has rank at most $C$.

Retain the ten entries $i\le 3-j$ of $X$ and the six strictly lower
entries of $(n_v^0)^{-1}L_v n_v^0$ at each place. This is injective
on the kernel: the infinitesimal symplectic adjoint determines the
other entries of $X$, and the flag tangent matrix in this basis is
strictly lower triangular. Hence the kernel has rank at most
$10C+6s\le16C$, with $s=\#\Phi.\iota$ in this general frame.
The sum of image and kernel ranks is at most $d+19C\le d+20C$.
Returning to the original pair by C65R1 proves the original
$\mathrm{RedTanA}$ bound on the span of its supported tangents.
% END C65R2

% BEGIN C65R3
\paragraph{C65R3. Dimension of the nonexceptional reducible loci.}
Let $R$ be any of the complete local deformation rings in
$\mathrm{RedA}$, representing lifts of the given residual pair $c$,
and let $P$ be its universal pair. Let $I_S$ be generated by the
off-diagonal entries of $P.r(g)$ of weight outside $S$. For any
of the eight rows other than $Y$ we claim
\[
                  \dim R/I_S\le d+20C+1 .
\]
If $I_S=R$, its closed locus is empty. Otherwise its residual
representation has support in $S$. Assume the prescribed type-$A$
witnesses: an entry of weight $-\sigma$ is nonzero, or entries of
weights $-\beta$ and $-\tau$ are both nonzero. Their weights must
belong to $S$ because the generating equations vanish after reduction.
Thus $S$ satisfies the criterion of C65R2.

Write $k$ for the residue field and $\epsilon^2=0$. A relative
derivation $\delta:R/I_S\to k$ at the closed point gives the
$\Lambda$-algebra map $u\mapsto\bar u+\epsilon\delta(u)$.
The resulting lift of $P$ has right logarithmic derivative
\[
   X(g)=\delta(P.r(g))c.r(g)^{-1},\qquad
   L_v=\delta(P.n_v)c.n_v^{-1},
\]
where the entries are first mapped to the quotient.
The coefficient of $\epsilon$ in each defining lift equation is
the corresponding tangent equation (UX.49); the support equations
hold already in $R/I_S$. Adic continuity here gives local constancy
because the maximal ideal of the dual numbers has square zero.
The map from derivations to these tangent vectors is injective.
Indeed zero derivatives on both matrices make the dual-number
lift agree with the constant lift. Universality gives equality of
the maps out of $R$, and surjectivity onto $R/I_S$ gives equality
on that ring. By C65R2 the relative cotangent dimension is at most
$d+20C$. Lifts of a basis give a surjection from a power-series
ring in at most that many variables over the complete DVR $\Lambda$,
by completeness and Nakayama's lemma. Its dimension bound, with the
one extra dimension of the DVR, proves $\mathrm{RedA}$.
% END C65R3

% BEGIN C65R4
\paragraph{C65R4. Curves avoiding the exceptional reducible component.}
Return to the original tower, space and ideal. Write $k=k_\Lambda$.
For a residual type-$A$ point satisfying either entry condition
of C65R3, take its local completion $f:A\to R$ representing lifts
(UX6). Write $\phi:B\to R$ for the composite and $U_v=T.U_1(v)$.
All the entries and ordinary characters below are their images on
these same rings. Let $d_v$ be the common local degree, so $s d_v=d$.
Put $N=2d+1-2s$. The source assumption $d\ge30(C+1)^2$ and
$s\le C$ give $s>0$, $d_v\ge30(C+1)$ and
\[
       d+20C+1<N,\qquad 2d+1+(2+4s)<N+d_v .
\]
Indeed $d>0$ and $s d_v=d$ force $s>0$, while the first
inequality follows from $d>20C+2s$ and the second from
$d_v>2+6s$.

Suppose minimal primes $P,Q$ of $R$ satisfy
$\dim R/(P+Q)\ge N+1$ and $J R\subset P$. Besides the eight
ideals $I_S$ for $S\ne Y$, use the finite family of ideals
\[
    B_{vij}=J R+(\varpi)+
        (\phi(\psi_{v,i}(h)-\psi_{v,j}(h)):h\in U_v),
           \qquad i\ne j ,
\]
where $\varpi$ is a uniformizer of $\Lambda$. By C65R3 each
$R/I_S$ in the first family has dimension less than $N$.
Take a prime $\mathfrak q$ containing a $B_{vij}$ but none of
these eight $I_S$, and set $\mathfrak b=\phi^{-1}\mathfrak q$.
The original finiteness condition on $J$ and coincidence on $U_v$
give $\dim B/\mathfrak b+d_v\le2d+1$ by (UX.55).
We claim that $\dim R/\mathfrak q\le\dim B/\mathfrak b+2+4s$
in both possible cases.

If $I_Y\not\subset\mathfrak q$, the generic representation avoids
all nine coordinate support patterns. The marked element has separate
eigenvalues, so its spectral projectors yield the coordinate lines.
The matrix-unit argument of (UX.46) gives absolute irreducibility,
and (UX.54) gives the claim. If $I_Y\subset\mathfrak q$, pass to
$\operatorname{Frac}(R/\mathfrak q)$. There the representation has
support in $Y$, and \emph{all four} weights of $Y$ occur. Indeed the
first four rows $0,1,2,3$ of (UX.46) omit from $Y$ respectively
$\tau,-\tau,-\beta,\beta$. If the indicated weight were absent,
the support would lie in the corresponding row, contrary to
$I_S\not\subset\mathfrak q$. Thus the exceptional trace-fibre
bound Y65P3 applies with all four nonvanishing hypotheses.
The two dimension inequalities give
$\dim R/\mathfrak q\le2d+1-d_v+2+4s<N$.

Cut $R/(P+Q)$ by $\varpi$. The principal ideal theorem leaves a
component of dimension at least $N$. Its prime contains neither an
$I_S$ from the eight nor a $B_{vij}$, by the two dimension tests.
Choose an element of each ideal outside this prime. Their product
stays nonzero on a one-dimensional local prime quotient of the
component, by the parameter argument of (UX.65). Thus there is a
prime $p\supset P+Q+(\varpi)$ with $\dim R/p=1$ avoiding both
families.

It remains essential to exclude $I_Y\subset p$ itself. Choose a
place $v$. The entry $n_{v,10}$ in $R$ is a unit because its residue
is nonzero by type $A$, so it stays nonzero in
$\operatorname{Frac}(R/p)$. If $I_Y\subset p$, Y65P1 would give
one fixed equality $\psi_{v,0}=\psi_{v,1}$ or
$\psi_{v,1}=\psi_{v,3}$ on all of $G_v$. Since $U_v\subset G_v$,
and $p$ already contains $JR+(\varpi)$, one of the avoided
$B_{vij}$ would lie in $p$. This contradiction excludes the ninth
pattern at the chosen curve.
% END C65R4

% BEGIN C65R5
\paragraph{C65R5. Every minimal component through a good rank-four point.}
The one-dimensional complete local domain $R/p$ from C65R4 embeds
in its normalization, a complete DVR of equal characteristic
with finite residue field $k'/k$. The resulting local map
$j:R\to k'[[t]]$ has kernel $p$ and respects the original
$\Lambda$-algebra structure, as in (UX.66). The point $jf$ sends
the base maximal ideal to series of zero constant term and kills
$J$. For every place and every distinct pair of indices, at least
one character difference on $U_v$ survives, by avoidance of all
$B_{vij}$. Its generic representation is absolutely irreducible
because C65R4 excluded all nine support patterns. The power-series
clause $\mathrm{p1}$ of the given potential modularity hypothesis
therefore supplies $\mathrm{LocZero}(J,jf)$. In particular
$Q\cap B=\phi^{-1}Q$ is a prime annihilated by this point on $B$,
since $Q\subset p$. The defining quantifier of $\mathrm{LocZero}$
gives $J\subset\phi^{-1}Q$. This propagates $JR\subset P$ to
$JR\subset Q$ across the chosen edge.

Now assume $\mathrm{SomeIn}(J,\ker x)$. Take one minimal prime of
$A$ witnessing this and any other minimal prime of $A$ under
$\ker x$. Both lift to minimal primes of $R$ by flatness of the
local completion; explicitly going down below its closed point lifts
either minimal prime, and then a minimal prime below that lift has the
same contraction. UX7 gives a presentation of $R$ by $b$ equations
in $a$ formal variables over $\Lambda$, where
$a-b\ge2d+2-2s=N+1$. The connectedness argument of (UX.62)
joins the two lifted minimal primes by a chain whose consecutive
intersections have dimension at least $N+1$. Successive application
of the preceding paragraph now puts $JR$ in the last prime.
Contracting gives $JA$ in the arbitrarily chosen minimal prime
of $A$.

Finally the original $\mathrm{Good4}$ condition consists of ordinary
goodness or one of the two type-$A$ alternatives of C65R3. In the
ordinary branch (UX.59) applies with the same fit, deep tower,
the supplied bounds and the two numerical inequalities of C65R4;
its use of the modularity hypothesis is through precisely the given
$\mathrm{fin}$ and $\mathrm{p1}$ fields. The preceding argument covers
both remaining branches. This proves the source
$\mathrm{ChainRankFour}$ implication
\[
       X.\mathrm{SomeIn}(J,\ker x)
          \quad\Longrightarrow\quad X.\mathrm{AllIn}(J,\ker x)
\]
under its stated cohomological and tower hypotheses and at the
unchanged threshold $30(C+1)^2\le d$. The construction of an
integral seed has a separate source statement.
% END C65R5
